\documentclass[10pt,a4paper]{amsart}
\usepackage[margin=19mm]{geometry}
\usepackage{amsmath,amssymb,mathtools}
\usepackage[scr=rsfs,cal=euler]{mathalfa}
\usepackage{xfrac}
\usepackage[unicode,hidelinks,bookmarks=false]{hyperref}
\newcommand{\ie}{i.e.\ }
\newcommand{\cf}{cf.\ }
\newcommand{\resp}{respectively\ }
\newcommand{\Subsection}{\subsection}
\providecommand{\nobreakdash}{\nobreak\hbox{-}}
\setlength{\emergencystretch}{2em}
\multlinegap0pt
\usepackage{amssymb}
\usepackage{enumerate}
\usepackage{dsfont}
\newtheorem{prop}{Proposition}
\newtheorem{lem}{Lemma}
\title{Joining rigidity for rational maps}
\author{Fabrizio Bianchi, Yan Mary He}
\date{Source: arXiv:2609.14890v1 (CC BY 4.0)}
\begin{document}
\maketitle

\noindent\emph{Source and licence.} Joining rigidity for rational maps. Version arXiv:2609.14890v1 (CC BY 4.0), distributed by arXiv under the Creative Commons Attribution 4.0 International licence, http://creativecommons.org/licenses/by/4.0/, as recorded in the arXiv licence metadata for this exact version and shown on the abstract page https://arxiv.org/abs/2609.14890v1. Full licence notice: \texttt{licenses/arxiv-2609.14890-CC-BY-4.0.txt}.\par\smallskip

\noindent\textit{Editorial scope.} Counted unit: the article's lem lem:return-equivar (source0.tex lines 2706--2722). The article's complete original proof of it is delivered with it (source0.tex lines 2724--2835, one \texttt{proof} environment of 112 lines). No mathematics is written, repaired or re-derived: the statement and the proof are the article's own lines, in the article's own notation, and no retained line is substituted or re-flowed.  Retained uncounted as the source-local context the proof needs: lines 640--648; lines 652--716.  Nothing was omitted from any retained span, so the package carries no omission record.  No retained line contains a citation, so the wrapper carries no bibliography and no bibliography entry.  No retained line contains a figure environment, an image-inclusion command or drawing code, so no image is delivered and the package declares no asset dependency.  Editorial formatting only, in addition: an AMS \texttt{amsart} preamble that restates the article's own declarations and macros used by the retained lines, namely the environments \texttt{prop}, \texttt{lem} on separate counters, together with the standard packages those lines need.  The retained environments print in this document as Proposition 1, Lemma 1 in order of appearance, whereas the article prints the same items with the numbering of its own full document.  The complete native source of the article as distributed in the arXiv e-print for this exact version is bundled unchanged as \texttt{sources/210369/source0.tex}.
\begin{equation}\label{eq:G}
G_{i,\hat z,n}
:=
\psi_{z_{-n}}
\circ
g_{i,\hat z,n}
\circ
\psi_{z_0}^{-1}.
\end{equation}

\begin{prop}
\label{prop:inverse branch-package}
For every $\varepsilon>0$ and $i\in\{1,2\}$, there exist a measurable
$\hat f_i$-invariant set
$\hat X_{i,\varepsilon}\subset\hat X_i$ with
$\hat\mu_i(\hat X_{i,\varepsilon})=1$ and a measurable function
$\rho_{i,\varepsilon}\colon
\hat X_{i,\varepsilon}\to(0,+\infty)$
such that the following hold for every
$\hat z\in\hat X_{i,\varepsilon}$ and every $n\geq1$.

\begin{enumerate}
\item
The inverse branch $g_{i,\hat z,n}$ is defined on
$B_{\mathrm{sph}}(z_0,\rho_{i,\varepsilon}(\hat z))$.

\item
The map $G_{i,\hat z,n}$ defined in \eqref{eq:G} is holomorphic and
univalent on
$D(0,\rho_{i,\varepsilon}(\hat z))$.

\item
For every $0<\lambda<1$, there exists a constant
$\kappa(\lambda)\geq1$, independent of $i$, such that, for every
$u,v\in
D(0,\lambda\rho_{i,\varepsilon}(\hat z))$,
we have
$$
\kappa(\lambda)^{-1}
\leq
\left|
\frac{G'_{i,\hat z,n}(u)}
     {G'_{i,\hat z,n}(v)}
\right|
\leq
\kappa(\lambda),
$$
and
$$
\kappa(\lambda)^{-1}|G'_{i,\hat z,n}(0)|
\leq
\frac{
\operatorname{diam}
G_{i,\hat z,n}
\bigl(D(0,\lambda\rho_{i,\varepsilon}(\hat z))\bigr)
}{
\lambda\rho_{i,\varepsilon}(\hat z)
}
\leq
\kappa(\lambda)|G'_{i,\hat z,n}(0)|.
$$

\item
For every $\hat z\in\hat X_{i,\varepsilon}$, we have
$
|G'_{i,\hat z,n}(0)|\to0
$
as $n\to\infty$. Consequently, for every $0<\lambda<1$,
$
\operatorname{diam}
G_{i,\hat z,n}
(D(0,\lambda\rho_{i,\varepsilon}(\hat z)))\to0.
$
\end{enumerate}
\end{prop}

\begin{lem}\label{lem:return-equivar}
Fix $z=(\hat x,\hat y)\in E^\sharp_{\mathrm{rec}}$, write
$Tz=(\hat x^+,\hat y^+)$, and set $\tau:=\tau(z)$. Then, after
restricting to sufficiently small neighbourhoods depending on $z$,
there exist a local biholomorphic germ $B_z$ and local
biholomorphic germs $A_{z,n}$ such that, for all sufficiently large
$n$, we have
\begin{equation}\label{eq:tilde-gauge}
\widetilde\Phi_{z,n+\tau}
=
A_{z,n}\circ
\widetilde\Phi_{Tz,n}\circ B_z.
\end{equation}
Moreover, every sequence $n_j\to\infty$ admits a subsequence along
which $A_{z,n_j}$ converges locally uniformly to a locally
biholomorphic germ.
\end{lem}

\begin{proof}
Since $Tz=\hat F^{-\tau}z$, we have
$\hat x^+=\hat f_1^{-\tau}\hat x$ and
$\hat y^+=\hat f_2^{-\tau}\hat y$. The inverse branches satisfy
$$
G_{1,\hat x,n+\tau}
=
G_{1,\hat x^+,n}\circ G_{1,\hat x,\tau}
\qquad
\text{and}
\qquad
G_{2,\hat y,n+\tau}
=
G_{2,\hat y^+,n}\circ G_{2,\hat y,\tau}.
$$
Recall that
$a_n(z)=G'_{2,\hat y,n}(0)/G'_{1,\hat x,n}(0)$. Differentiating the
preceding identities gives
$$
a_{n+\tau}(z)
=
a_n(Tz)\lambda_z
\qquad
\text{and}
\qquad
\lambda_z
:=
\frac{G'_{2,\hat y,\tau}(0)}
     {G'_{1,\hat x,\tau}(0)}
\in\mathbb C^*.
$$
Using the definitions of the normalized transfer maps and the cocycle
relations, we obtain
$$
\Phi_{z,n+\tau}
=
L_{z,n}\circ
\Phi_{Tz,n}\circ
G_{1,\hat x,\tau},
\quad \text{where}
\quad
L_{z,n}
:=
G_{2,\hat y,\tau}^{-1}
\circ
G_{2,\hat y^+,n}^{-1}
\circ
(\lambda_z\cdot\operatorname{id})
\circ
G_{2,\hat y^+,n}.
$$
All these identities are understood on sufficiently small
neighbourhoods on which the compositions are defined.
Passing to the fixed reference charts gives
\eqref{eq:tilde-gauge}, with
$$
A_{z,n}
:=
k_{\hat y}\circ L_{z,n}\circ k_{\hat y^+}^{-1}
\qquad
\text{and}
\qquad
B_z
:=
h_{\hat x^+}^{-1}\circ
G_{1,\hat x,\tau}\circ h_{\hat x}.
$$
The germ $B_z$ is a local biholomorphism depending only on the finite
return segment.
For fixed $z$, the only $n$-dependent part of $A_{z,n}$ is 
$$
H_n
:=
G_{2,\hat y^+,n}^{-1}
\circ
(\lambda_z\cdot\operatorname{id})
\circ
G_{2,\hat y^+,n}.
$$
By Proposition \ref{prop:inverse branch-package}, the maps
$G_{2,\hat y^+,n}$ are univalent on a common disc and satisfy uniform
Koebe estimates there. Choose $t>0$ sufficiently small so that
$G_{2,\hat y,\tau}^{-1}$ is defined on $D(0,t)$. Since $\lambda_z$
is fixed, we may then choose $0<s<t$, depending on $z$ but not on
$n$, such that
$$
\lambda_z G_{2,\hat y^+,n}(D(0,s))
\subset
G_{2,\hat y^+,n}(D(0,t))
$$
for every $n$. Indeed, the Koebe distortion and growth estimates give
uniform inner and outer bounds for these images in terms of
$|G'_{2,\hat y^+,n}(0)|$.
It follows that $H_n$ is defined on $D(0,s)$ for every $n$ and
satisfies
$H_n(D(0,s))\subset D(0,t)$.
Hence $\{H_n\}$ is a normal family on $D(0,s)$.
Moreover,
$H_n'(0)=\lambda_z$
for every $n$.
Consequently, every locally uniform cluster limit of the $H_n$ has
non-zero derivative at the origin and is therefore locally
biholomorphic after restriction.

Since $A_{z,n}$ is obtained from $H_n$ by pre- and
post-composition with fixed local biholomorphisms,
after restricting
to a sufficiently small fixed neighbourhood of the base point, the
normality of $\{H_n\}$ gives normality of
$\{A_{z,n}\}$. Every cluster limit is locally biholomorphic as
the same holds for the corresponding cluster limit of $\{H_n\}$.
\end{proof}

\end{document}
